\section{Sobre el Comité}

\subsection{Historia y Naturaleza del Comité}

El Fondo de las Naciones Unidas para la Infancia (UNICEF) es la agencia de las Naciones Unidas especializada en la infancia. Este órgano focaliza sus debates e investigaciones en temas como la educación, la salud y la nutrición infantil, la protección contra la violencia, la explotación y frente al cambio climático así como promueve el salvaguardo de los derechos de la infancia.\footnote{UNICEF. (s. f.). ¿Qué hacemos? \url{https://www.unicef.org/es/que-hacemos}} UNICEF lleva a cabo todos sus proyectos con una equidad fundamental y opera en un contexto tanto global como específico buscando salvaguardar y promover el bienestar de las generaciones más jóvenes de la sociedad.

Este organismo fue creado el 11 de diciembre de 1946 por la Asamblea General de las Naciones Unidas basado en el documento fundamental de la Convención sobre los Derechos del Niño, adoptada en 1989 y convirtiéndose en el documento fundamental de derechos más ratificado del mundo.\footnote{UNICEF España. (s. f.). 60 años trabajando para la infancia. \url{https://www.unicef.es/noticia/60-anos-trabajando-para-la-infancia}} \footnote{Oficina del Alto Comisionado de las Naciones Unidas para los Derechos Humanos. (s. f.). Convención sobre los Derechos del Niño. \url{https://www.ohchr.org/en/instruments-mechanisms/instruments/convention-rights-child}}

\subsection{Funciones}

En sus estados más tempranos, UNICEF era un fondo temporal para la protección y recuperación de los niños durante el periodo de posguerra tras la Segunda Guerra Mundial y se convirtió en un organismo fijo de las Naciones Unidas en el año 1953.\footnote{UNICEF España. (s. f.). 60 años trabajando para la infancia. \url{https://www.unicef.es/noticia/60-anos-trabajando-para-la-infancia}} Tras la Segunda Guerra Mundial, los niños y niñas del mundo se encontraron bajo peligros inminentes en cuanto a salud, bienestar y educación, por lo que este órgano fue creado. UNICEF lidia con múltiples complicaciones a diario ya sean obstáculos físicos en zonas de conflicto u otros que dificultan la protección de los derechos de los menores, pero aun así, los logros de esta organización son incontables.

UNICEF a día de hoy es el foro fundamental para la cooperación de los países miembros, agencias especializadas en la infancia y actores externos a la hora de tratar temas que conciernen a todos los niños y niñas tanto de sus respectivas comunidades como en la escena global.\footnote{UNICEF. (s. f.). ¿Qué hacemos? \url{https://www.unicef.org/es/que-hacemos}} A lo largo de los años, UNICEF se ha convertido en uno de los pilares de las Naciones Unidas bajo marcos cambiantes que se adaptan a las circunstancias de actualidad en cada momento. La volatilidad del escenario global en cuanto a crisis humanitarias, faltas a los derechos del niño y el efecto que estos eventos causan en los infantes de la sociedad son objetivos claros que UNICEF trabaja día a día de cara a erradicar y cumplir con los objetivos impuestos en los Objetivos de Desarrollo Sostenible para 2030.\footnote{Organización de las Naciones Unidas. (s. f.). El sistema de las Naciones Unidas. \url{https://www.un.org/es/about-us/un-system}}

Actualmente la base de UNICEF está establecida principalmente en Nueva York desde donde coordina operaciones en más de 190 países distintos. Bajo el objetivo principal de proteger y promover el bienestar y la salud de los más pequeños, UNICEF ha llevado a cabo miles de misiones humanitarias y proyectos de desarrollo anuales en los países miembros. A cargo de los niños y niñas de todo el mundo, UNICEF tiene la capacidad de llevar a cabo

\subsection{Logros}

Entre los proyectos de mayor impacto positivo que el organismo ha protagonizado se encuentran su liderazgo ante la inmunización masiva, proporcionando vacunas a alrededor de un 50\% de la población global así como su actuación como distribuidor de Alimentos Terapéuticos Listos para el Consumo (ATLCs). Además, entre otros se encuentra WASH, que es un mecanismo utilizado en zonas vulnerables o en estado de crisis para proporcionar sistemas de agua mediante energía solar, habilitando agua potable a las poblaciones de todo el mundo.\footnote{Forbes. (2025, 24 de febrero). 10 maneras en que la inversión en el bienestar infantil cambió el mundo. \url{https://www.forbes.com/sites/unicefusa/2025/02/24/10-ways-investment-in-childrens-well-being-changed-t} he-world/}

Específicamente, el Subsidio de Apoyo a la infancia en Tailandia (Child Support Grant) en 2015 protege a día de hoy a millones de niños en vulnerabilidad alta de circunstancias de pobreza extrema, convirtiéndose así en una política pública de forma permanente a nivel nacional. Este proyecto duplicó las inversiones en fondos públicos de Tailandia. Otros proyectos incluyen la construcción de infraestructura principalmente dirigida al acceso a la educación en Yemen o el proyecto de Salud Preventiva en Kosovo en 2024.\footnote{UNICEF España. (s. f.). Memoria 2023. \url{https://www.unicef.es/memoria/2023}}

Este organismo se basa en procedimientos similares a aquellos de las otras cinco comisiones que conforman asamblea general junto con la UNICEF. A pesar de la falta de poder de ejecución, UNICEF ha cambiado millones de vidas a lo largo de los años y continúa siendo el núcleo de la protección infantil como uno de los seis órganos principales de la Asamblea General.\footnote{Organización de las Naciones Unidas. (s. f.). El sistema de las Naciones Unidas. \url{https://www.un.org/es/about-us/un-system}}

\subsection{Retos Actuales}

Aún teniendo en cuenta los numerosos éxitos del organismo, UNICEF se enfrenta a ciertos obstáculos recurrentes a la hora de cumplir con su misión. En el siglo XXI, UNICEF hace frente a problemas de naturaleza híbrida y volátil dada la interdependencia entre naciones tanto comercial como educativa, además de climática y financiera. A la hora de proveer ayuda humanitaria y asistencia de desarrollo para los niños necesitados, este órgano se enfrenta a una grave dependencia en fondos externos y al desarrollo de los propios conflictos que obstaculizan una actuación rápida y eficaz.\footnote{Noticias ONU. (2025, enero). [Noticia sobre los derechos de la infancia]. \url{https://news.un.org/es/story/2025/01/1535651}} No obstante, a pesar de varios resultados positivos, UNICEF continúa luchando contra una desigualdad creciente, el cambio climático y avances tecnológicos de rapidez incalculable que crean una incertidumbre monumental a la hora de mitigar el impacto de dichos factores en la infancia.\footnote{World Economic Forum. (2025, enero). Avances tecnológicos y desarrollo humano: Una historia de dos mundos. \url{https://es.weforum.org/stories/2025/01/avances-tecnologicos-y-desarrollo-humano-una-historia-de-dos-m} undos/}

Pese a una lucha continua contra las crisis humanitarias, especialmente de cara a los infantes, UNICEF continúa enfrentándose a un escenario geopolítico de conflicto creciente. Infraestructuras e instituciones cuyo objetivo es la protección de los niños y niñas se han convertido en un objetivo claro a la hora de atacar a un estado, pasando así la legislatura internacional a un segundo plano de manera cada vez más frecuente.

En cuanto a la estructura financiera, cada vez sufre de mayor inestabilidad, especialmente en países en vías de desarrollo, además de crecientes deudas soberanas en los países que cada vez dificultan más la inversión en la infancia.\footnote{UNICEF. (s. f.). Cómo mantener a los niños seguros en Internet. \url{https://www.unicef.org/protection/keeping-children-safe-online}}

Por último, el fácil y desordenado acceso a las plataformas digitales suponen uno de los mayores retos para el organismo dado el peligro que presentan de cara a la seguridad de todos los niños y niñas del mundo. Un entorno digital velozmente cambiante nos adentra en una era de incertidumbre, creando así nuevos peligros que se adentran en todos los ámbitos de las vidas de todos los niños que, en conjunto con crecientes desigualdades, promueven un acceso indebido a datos que deberían permanecer bajo una privacidad estricta.\footnote{UNICEF España. (s. f.). Nuevo informe sobre el impacto de la tecnología en la infancia y la adolescencia. \url{https://www.unicef.es/noticia/nuevo-informe-sobre-el-impacto-de-la-tecnologia-en-la-infancia-y-la-adolesc} encia}

Es por esto que UNICEF continúa luchando contra todas las adversidades para promover el cambio y el desarrollo positivo, especialmente en aquellas regiones que más lo necesitan.
